\documentclass[11pt,reqno]{amsart}
\usepackage[T1]{fontenc}
\usepackage[utf8]{inputenc}
\usepackage{lmodern}
\usepackage[a4paper,margin=27mm]{geometry}
\usepackage{amsmath,amssymb,mathtools,mathrsfs}
\usepackage{microtype}
\usepackage{xurl}
\usepackage[hidelinks]{hyperref}
\hypersetup{pdftitle={Weyl growth and tracial limits for arithmetic translations},pdfauthor={R. Arun Chandru}}
\numberwithin{equation}{section}
\setcounter{tocdepth}{1}
\allowdisplaybreaks[2]
\emergencystretch=2em
\raggedbottom
\newcommand{\R}{\mathbb R}
\newcommand{\C}{\mathbb C}
\newcommand{\E}{\mathbb E}
\newcommand{\Tr}{\operatorname{Tr}}
\newcommand{\tr}{\operatorname{tr}}
\newcommand{\supp}{\operatorname{supp}}
\newcommand{\dd}{\,\mathrm d}
\newcommand{\ii}{\mathrm i}
\newcommand{\one}{\mathbf 1}
\newcommand{\norm}[1]{\lVert#1\rVert}
\newcommand{\abs}[1]{\lvert#1\rvert}
\newcommand{\ip}[2]{\langle#1,#2\rangle}
\title[Weyl growth and arithmetic tracial limits]{Weyl growth and tracial limits\texorpdfstring{\\}{ }for arithmetic translations}
\author{R. Arun Chandru}
\email{arun.chandru@gmail.com}
\date{27 September 2026}
\subjclass[2020]{Primary 35P20, 11F30; Secondary 47B35, 60B20, 60F10}
\keywords{Weyl coefficient, holomorphic cusp form, killed translation,
Rankin--Selberg prime number theorem, Toeplitz matrix, tracial spectral law,
moderate deviations, autocorrelation}
\begin{document}
\begin{abstract}
We determine the growing-window tracial pressure of the actual
prime-power translation operator attached to a fixed holomorphic
cuspidal newform of trivial character. Every intermediate translation
is killed at the interval boundary. If the interval has length $\ell$,
the logarithmic tracial exponential at fixed coupling $s$ is
$s^2\ell^2/12+O_{f,s}(\ell\sqrt{\log\ell})$, with a bounded lower
error. The corresponding logarithmic Weil operator has Weyl
coefficient $K_{f,\ell}$ satisfying
$\log K_{f,\ell}=\ell^2/48+\log(2\ell/\sqrt N)
+O_f(\ell\sqrt{\log\ell})$.
A uniform coupling estimate gives the spectral large-deviation law
at every diverging intermediate scale up to $\ell^2$.
At the central scale $\ell$, the arithmetic trace converges to the
$\sqrt a$ variance-profile Toeplitz law, whose second and fourth
moments are $1/3$ and $4/15$. We determine its logarithmic tail
coefficient, equal to $-3$ in both directions. A sharp weighted
autocorrelation identity explains the pressure constant. We prove
global $L^2$ stability modulo affine phase, with optimal square-root
exponent and sharp small-deficit squared-distance coefficient $6$.
An explicit taper family preserves total leading prime variance
while changing the boundary pressure. Higher prime powers and ramification change
the logarithmic pressure by a bounded amount, without being small
in operator norm. Fixed tuples of inequivalent forms give an
orthogonally invariant, noncommuting colored limit. The proofs
retain boundary survival through a uniform finite-grid comparison;
fixed-moment convergence alone is not used to exchange the window
and coupling limits.
\end{abstract}

\maketitle
\tableofcontents

\section{Introduction}\label{sec:introduction}

For an operator whose principal growth is logarithmic in frequency,
a bounded perturbation can change the leading Weyl coefficient.
The coefficient records an exponential statistic of the perturbation,
not merely its average. When the perturbation is a sum of translations
on an interval, the relevant average also remembers the boundary:
a path contributes only if every one of its intermediate vertices
remains inside the interval. Arithmetic translations make this
geometry particularly rigid. Their lengths are logarithms of prime
powers, and unique factorization determines exactly which words can
return to their starting point.

This paper studies that coefficient as the interval grows. The
translation alphabet then grows exponentially, its entries have
arithmetic signs, and a fixed-alphabet Weyl formula does not estimate
the resulting coefficient. Our main result is a sharp pressure law
for this growing arithmetic operator. A second result identifies its
central tracial distribution with a known variance-profile Toeplitz
law and determines the sharp logarithmic tails of that law. These
are two different limiting questions. We give a uniform estimate
which also fills the intermediate spectral scales.

\subsection*{The arithmetic operator and the main estimates}

Fix a primitive holomorphic cuspidal newform $f$ of even weight
$2\kappa\ge2$, exact level $N$, and trivial character. We use its
real, unitary-normalized Hecke coefficients. Write $\alpha_p,\beta_p$
for its good local parameters, so
$\alpha_p\beta_p=1$, $\abs{\alpha_p}=\abs{\beta_p}=1$, and
$\lambda_f(p)=\alpha_p+\beta_p$.
For $I_\ell=(0,\ell)$ let
\[
 (S_hu)(x)=\one_{I_\ell}(x-h)u(x-h),\qquad x\in I_\ell,
\]
where $u$ is extended by zero. In particular $S_h^*=S_{-h}$.
The operator under study is the finite selfadjoint sum
\begin{equation}\label{eq:W-intro}
 W_{f,\ell}=\sum_{p^m<e^\ell}b_f(p^m)
             (S_{m\log p}+S_{-m\log p}).
\end{equation}
At good primes,
$b_f(p^m)=(\log p)p^{-m/2}(\alpha_p^m+\beta_p^m)$;
the actual ramified coefficients are specified in
Section~\ref{sec:objects}. These are logarithmic-derivative
coefficients, not the ordinary Hecke coefficients at $p^m$.

The translation algebra has a normalized positive trace $\tau_\ell$.
It is not the Hilbert-space trace. If a word has successive
displacements $a_0=0,a_1,\ldots,a_r$, its trace is
\begin{equation}\label{eq:word-intro}
 \one_{\{a_r=0\}}
 \left(1-\frac{\max_j a_j-\min_j a_j}{\ell}\right)_+.
\end{equation}
We construct this trace on the actual operators below. Put
\[
 P_{f,\ell}(s)=\log\tau_\ell(e^{sW_{f,\ell}}).
\]
Our uniform estimate is the following precisely quantified assertion:
for each fixed $f$ and $C>0$, there are constants $A_{f,C}$ and
$\ell_0(f,C)$ such that, whenever $\ell\ge\ell_0(f,C)$ and
$2\le\abs t\le C\ell$,
\begin{equation}\label{eq:uniform-main}
 \frac{t^2}{12}-A_{f,C}
 \le P_{f,\ell}(t/\ell)
 \le\frac{t^2}{12}
        +A_{f,C}\bigl(1+\abs t\sqrt{\log(2+\abs t)}\bigr).
\end{equation}
The form and $C$ are fixed; the estimate is uniform in the displayed
range of $t$ and $\ell$. There is no assertion of a threshold uniform
over varying levels or weights. In particular, for every fixed real $s$,
\begin{equation}\label{eq:pressure-main}
 \lim_{\ell\to\infty}\ell^{-2}P_{f,\ell}(s)=\frac{s^2}{12}.
\end{equation}
For $s\ne0$ its lower error is bounded and its upper error is
$O_{f,s}(\ell\sqrt{\log(2+\ell)})$; $s=0$ is exact.

The exterior-zero logarithmic Weil operator associated with this
normalization has fixed-window counting function
$N_{f,\ell}(E)\sim K_{f,\ell}e^{E/2}$, where
\begin{equation}\label{eq:K-intro}
 K_{f,\ell}=\frac{2\ell}{\sqrt N}\tau_\ell(e^{W_{f,\ell}/2}).
\end{equation}
We recall the fixed-window input and verify its normalization in
Section~\ref{sec:objects}. Thus
\begin{equation}\label{eq:K-main}
 \begin{split}
 \frac{\ell^2}{48}+\log\frac{2\ell}{\sqrt N}-O_f(1)
 &\le\log K_{f,\ell}\le
 \frac{\ell^2}{48}+\log\frac{2\ell}{\sqrt N}
       +O_f(\ell\sqrt{\log(2+\ell)}).
 \end{split}
\end{equation}
The high-energy limit defining $K_{f,\ell}$ is taken first, at a
fixed window. The subsequent limit in~\eqref{eq:K-main} does not
give a window-uniform onset for the counting asymptotic.

\subsection*{Central, intermediate, and exponentially weighted scales}

The probability spectral laws of $W_{f,\ell}/\ell$ in $\tau_\ell$
converge, with all moments and all real exponential moments, to one
symmetric measure $\nu$, independent of $f$. It is the expected
limiting spectral law of the Hermitian Gaussian Toeplitz matrices
\begin{equation}\label{eq:toeplitz-intro}
 T_n=\sum_{j=1}^{n-1}\frac{\sqrt j}{n}
       (z_jS_j+\overline z_jS_j^*),
 \qquad \E\abs{z_j}^2=1,\quad \E z_j^2=0,
\end{equation}
where the $z_j$ are independent circular Gaussians and the shifts
act on $\C^n$. Its first nonzero moments and sharp tail rates are
\begin{equation}\label{eq:central-intro}
 \int x^2\dd\nu=\frac13,\qquad
 \int x^4\dd\nu=\frac4{15},\qquad
 \lim_{R\to\infty}\frac{\log\nu([R,\infty))}{R^2}
 =\lim_{R\to\infty}\frac{\log\nu(( -\infty,-R])}{R^2}=-3.
\end{equation}
This is not a Gaussian distribution: a Gaussian of variance $1/3$
has fourth moment $1/3$.

Between this central scale and the scale $\ell^2$, choose any positive
$r_\ell\to\infty$ with $r_\ell=O(\ell)$. The spectral laws of
$W_{f,\ell}/(\ell r_\ell)$ satisfy a large-deviation principle with
speed $r_\ell^2$ and rate $3x^2$. No lower rate of divergence for
$r_\ell$ is imposed. At the endpoint $r_\ell=\ell$ this determines
the rare tracial spectral mass carrying the Weyl coefficient.
For the central limiting measure, a separate positive-word comparison
proves
\begin{equation}\label{eq:central-mgf-intro}
 \frac{s^2}{12}\le
 \log\int e^{sx}\dd\nu(x)
 \le\frac{s^2}{12}+O\bigl(\abs s\sqrt{\log(2+\abs s)}\bigr)
 \quad (\abs s\to\infty).
\end{equation}
This last statement is not obtained by interchanging a fixed-moment
limit with a growing coupling.

There is a useful robustness statement beyond the untapered operator.
For $g_\varepsilon(a)=\sqrt{1+\varepsilon\cos(2\pi a)}$,
with $-9/20\le\varepsilon\le24/(4\pi^2-3)$, multiply each
coefficient at $p^m$ by $g_\varepsilon(m\log p/\ell)$.
The total leading first-prime variance remains $\ell^2/2$,
but the pressure constant becomes $1/12-\varepsilon/(4\pi^2)$.
Section~\ref{sec:tapers} derives this from a general boundary-profile
criterion. This is an explicitly defined deformation, not an
alternative formula for the untapered Weil operator.

\subsection*{Why the boundary changes the constant}

If every surviving fraction in~\eqref{eq:word-intro} were replaced by
one, primewise balance would give an independent-circle product with
logarithm $s^2\ell^2/2+O_{f,s}(1)$. The boundary reduces its leading
coefficient from $1/2$ to $1/12$. The reduction is determined by the
weighted autocorrelation problem
\begin{equation}\label{eq:variational-intro}
 \sup_{\norm u_{L^2(0,1)}=1}
 \int_0^1 a\left|\int_0^{1-a}\overline{u(x)}u(x+a)\dd x\right|^2
 \dd a=\frac1{12}.
\end{equation}
We give an exact nonnegative deficit identity and show that the
maximizers are precisely the flat-modulus linear-phase profiles.
More strongly, if $D(u)=1/12-J(u)$ for a unit profile and $d(u)$
is its $L^2$ distance to that family, then
\begin{equation}\label{eq:stability-intro}
 d(u)^2\le C D(u),\qquad
 d(u)^2\le6D(u)+O(D(u)^{3/2})\quad\text{as }D(u)\downarrow0.
\end{equation}
The global constant is absolute but not numerically specified;
the local coefficient $6$ is optimal. The proof combines weak
quartic splitting modulo modulation with an explicit transverse
Hessian calculation. It requires more than the weak modulus
control visible in the deficit identity alone.
A separate finite-dimensional deficit identity controls the
variance of every vector of the Gaussian comparison matrix,
not just a proposed test vector.

Two further distinctions are essential. First, powers of one prime
share one phase: they are not independent random coefficients.
Their complete contribution can nevertheless be restored at bounded
logarithmic cost by a relative tracial Taylor estimate. Second,
rounding translation lengths can create false returns. We keep an
independent phase for each prime while rounding, so that these aliases
are removed exactly, and enlarge the interval to dominate every
survival mask multiplicatively. An additive error in each word would
not suffice at exponential scale.

\subsection*{Relation to prior work and scope of the contribution}

Laptev and Weth~\cite{LW} treat the exterior-zero logarithmic
Laplacian; the bounded translation-algebra extension used
in~\eqref{eq:K-intro} is established in~\cite{ChandruW}.
Here the new problem is the growth of its actual arithmetic
coefficient as the window and alphabet grow. The arithmetic interval
operator and its critical heat residue also appear in~\cite{ChandruB}.
Our proof uses
Deligne's local bound~\cite{Deligne} and the fixed-form
Rankin--Selberg prime number theorem in the form of
Gillespie and Ji~\cite{GJ}.

Bryc, Dembo and Jiang~\cite{BDJ} obtain the classical scalar law,
and Bose, Saha and Sen~\cite{BSS} explicitly allow variance profiles.
Their theorem contains the scalar law in~\eqref{eq:toeplitz-intro}
as the profile $\sigma(a)=\sqrt a$. Colored pairing limits and the
failure of asymptotic freeness for Toeplitz matrices are already
studied by Bose, Hazra and Saha~\cite{BHS}. We prove the occurrence
of these structures in the original arithmetic translation trace,
including every local factor, and give the profile-specific pressure
and tail calculation. 

The unweighted problem of Garsia, Rodemich and Rumsey~\cite{GRR}
appears in the Toeplitz norm analysis of Sen and Vir\'ag~\cite{SV};
Onatski~\cite{Onatski} studies related rectangular geometries.
Manov~\cite{Manov} treats weighted linear functionals on positive
definite functions. Our quadratic functional retains both the lag
weight and the square-autocorrelation constraint. The direct linear
relaxation does not give its sharp constant, as explained in
Section~\ref{sec:profile}. 

The Gaussian tools are standard concentration and matrix-series
estimates~\cite{Gross,Ledoux,Tropp}. The specific analytic work is a
uniform comparison that retains killed paths, controls the actual
prime variances in moving bins, and permits a growing exponential
coupling. The large-deviation conclusions are consequences of that
pressure estimate by exponential tilting, whose proof is included.
All statements are unconditional for the fixed holomorphic forms
specified above, including forms with complex multiplication.

\section{The interval trace and the Weyl coefficient}\label{sec:objects}

\subsection*{A trace on the concrete translation algebra}

Let $\mathcal A_\ell$ be the unital norm-closed algebra on
$L^2(I_\ell)$ generated by the $S_h$ and bounded multiplication
operators. Products before closure have the form
$M_gS_h$ with the appropriate overlap indicators absorbed in $g$.
The definition does not require a unique representation as a finite sum.
For $v_t(x)=\ell^{-1/2}e^{\ii tx}$ define on finite sums
\begin{equation}\label{eq:trace-construction}
 \tau_\ell(B)=\lim_{R\to\infty}\frac1{2R}
                 \int_{-R}^R\ip{v_t}{Bv_t}\dd t.
\end{equation}
The inner product is conjugate-linear in its first argument.
For a monomial the integrand is $e^{-\ii th}$ times its surviving
coefficient integral divided by $\ell$. Thus the limit is zero
if $h\ne0$, and is $\ell^{-1}\int g$ if $h=0$.
Each approximating functional is positive, unital, and contractive.
The limit therefore extends continuously to $\mathcal A_\ell$.
For two monomials the traces of their products vanish unless the
net displacement is zero; in that case a translation of the
integration variable proves cyclicity. Linearity and continuity
give a tracial state on the full algebra.

If $a_j=h_1+\cdots+h_j$, then
$S_{h_1}\cdots S_{h_r}$ acts at $x$ only when all $x-a_j$ lie in
$I_\ell$. For a closed word this set of starting positions has
length $(\ell-\max a_j+\min a_j)_+$. This proves
\eqref{eq:word-intro}. The formula keeps the boundary at each
step; it is generally false that $S_hS_k=S_{h+k}$.

For a bounded selfadjoint $B\in\mathcal A_\ell$, the restriction of
$\tau_\ell$ to its continuous functional calculus determines a
probability measure. We use that measure to interpret spectral
indicators. Equivalently one may pass to the finite von Neumann
algebra of the tracial GNS representation, after dividing out the
trace-null ideal. Its normal trace defines $\tau_\ell(\one_E(B))$
for every Borel set $E$. No trace on all of
$\mathcal B(L^2(I_\ell))$ is asserted.

\subsection*{Actual local factors}

At $p\nmid N$ let the unitary Euler factor be
$((1-\alpha_pp^{-s})(1-\beta_pp^{-s}))^{-1}$.
If $p\Vert N$, it is $(1-\epsilon_pp^{-s-1/2})^{-1}$ with
$\epsilon_p\in\{1,-1\}$; for trivial character it has degree
zero when $p^2\mid N$; these are the classical newform local
factors~\cite{AL}. Consequently
\begin{equation}\label{eq:local-coefficients}
 b_f(p^m)=
 \begin{cases}
  (\log p)p^{-m/2}(\alpha_p^m+\beta_p^m),&p\nmid N,\\
  \epsilon_p^m(\log p)p^{-m},&p\Vert N,\\
  0,&p^2\mid N.
 \end{cases}
\end{equation}
The operator~\eqref{eq:W-intro} is bounded at each fixed $\ell$
because it is a finite sum. Terms with $p^m=e^\ell$ vanish as
operators, so either endpoint convention gives the same $W_{f,\ell}$.
We write $W_{f,\ell}^{(1)}$ for its good first-prime part and put
\begin{equation}\label{eq:first-coeff}
 b_p=\lambda_f(p)\frac{\log p}{\sqrt p},\qquad h_p=\log p,
 \qquad W^{(1)}=\sum_{p\nmid N,\ h_p<\ell}b_p(S_{h_p}+S_{-h_p}).
\end{equation}
In expressions involving $b_p$, unqualified prime sums run
over primes good for $f$.

\subsection*{The fixed-window input and its normalization}

Use $\widehat{u^0}(\xi)=\int_\R e^{-\ii x\xi}u^0(x)\dd x$.
The exterior-zero logarithmic operator $H_{I_\ell}$ is defined by
the closed semibounded form with symbol $\log\abs\xi$ and
Fourier factor $(2\pi)^{-1}$. Its negative part is bounded because
$\abs{\widehat{u^0}(\xi)}\le\sqrt\ell\norm u_2$ and
$\log\abs\xi$ is integrable near zero. Its form domain is the
space of supported $L^2$ functions with finite positive logarithmic
Fourier energy. This is not the spectral logarithm of a Dirichlet
Laplacian.

The fixed-window result we use from~\cite{ChandruW}, in one
dimension, states that for $d>0$, $c\in\R$, bounded selfadjoint
$B\in\mathcal A_\ell$, and compact selfadjoint $R$,
\begin{equation}\label{eq:fixed-weyl}
 N_{dH_{I_\ell}+c+B+R}(E)
 \sim\frac\ell\pi e^{-c/d}\tau_\ell(e^{-B/d})e^{E/d}
 \qquad(E\to\infty).
\end{equation}
All the data in this statement are fixed while $E$ tends to
infinity. This external fixed-window input, like the arithmetic
prime number theorem used below, is separated from the new
growing-window estimate.

Let $\psi=\Gamma'/\Gamma$. Define $A_{f,\ell}$ on the same
logarithmic form domain by the Fourier symbol
\[
 \log N-2\log(2\pi)+2\Re\psi(\kappa+\ii\xi)
\]
and the bounded perturbation $-W_{f,\ell}$. Since
\[
 r_\kappa(\xi)=2\Re\psi(\kappa+\ii\xi)-2\log\abs\xi
       \in L^1(\R),
\]
its compressed multiplier $R_{\kappa,\ell}$ is trace class.
Indeed, factor it as $T^*M_{\operatorname{sgn}r_\kappa}T$,
where $Tu=(2\pi)^{-1/2}\abs{r_\kappa}^{1/2}\widehat{u^0}$;
then $\norm T_{\mathrm{HS}}^2=\ell\norm{r_\kappa}_1/(2\pi)$.
The integrability follows from the digamma asymptotic at infinity
and the integrable logarithm at zero~\cite{DLMF}. Thus, as an
identity of closed forms and their operators,
\begin{equation}\label{eq:gamma-identity}
 A_{f,\ell}=2H_{I_\ell}
       +\bigl(\log N-2\log(2\pi)\bigr)I
       -W_{f,\ell}+R_{\kappa,\ell}.
\end{equation}
Using $d=2$ in~\eqref{eq:fixed-weyl} gives exactly
\eqref{eq:K-intro}. No uniform bound on
$\norm{R_{\kappa,\ell}}_1$ as the window grows is needed for
this fixed-window identification.

\section{Arithmetic variance and bounded local restoration}\label{sec:arithmetic}

\subsection*{The quantitative prime input}

For our fixed form the Rankin--Selberg prime number theorem gives
\begin{equation}\label{eq:RS}
 A_f(X):=\sum_{\substack{p\le X\\p\nmid N}}
              \lambda_f(p)^2\log p
       =X+O_f\bigl(Xe^{-c_f\sqrt{\log X}}\bigr),\qquad c_f>0.
\end{equation}
One may use~\cite[Lemma 3.2]{GJ} over $\mathbb Q$. The representation
is self-contragredient. That statement first includes prime powers;
Deligne's bound removes their good contribution with
$O(\sqrt X\log^2(2X))$ error, and the finitely many bad primes
are harmless at fixed $f$. Decreasing $c_f$ absorbs both errors.
No conjectural prime-splitting assertion is needed.

Partial summation with $(\log X)/X$ gives
\begin{equation}\label{eq:variance}
 V_f(H):=\sum_{\substack{h_p<H\\p\nmid N}}b_p^2
       =\frac{H^2}{2}+O_f(1),\qquad H\ge0.
\end{equation}
The bounded error, rather than just $o(H^2)$, follows because
$(1+H)e^{-c_f\sqrt H}$ is integrable on the positive half-line.
Endpoint atoms are covered by $b_p^2\le4(\log p)^2/p$.
The same calculation on a logarithmic interval $[H,H+\Delta)$,
with $0\le H<H+\Delta\le\ell$, gives the conservative bound
\begin{equation}\label{eq:bin}
 \sum_{H\le h_p<H+\Delta}b_p^2
 =\int_H^{H+\Delta}v\dd v
   +O_f\bigl((\ell+1)^2e^{-c_f\sqrt H}\bigr),
\end{equation}
after decreasing $c_f$ if necessary and treating the finite initial
interval by its constant. This form is valid even when $\Delta>1$.
It follows by applying Stieltjes partial summation at the two
endpoints and integrating the error over the interval; the bound
$(\ell+1)(1+\Delta)\le(\ell+1)^2$ suffices.

Two further consequences used repeatedly are
\begin{equation}\label{eq:variance-measure}
 \ell^{-2}\sum_{h_p<\ell}b_p^2\delta_{h_p/\ell}
       \Longrightarrow a\dd a\quad\hbox{on }[0,1],
 \qquad \sum_p b_p^4<\infty.
\end{equation}
The fourth-power sum has an absolute majorant independent of $f$;
also $\abs{b_p}\le2(\log p)/\sqrt p\le4/e$.

\subsection*{A relative tracial Taylor estimate}

We record the noncommutative estimate that will justify changing
local variables without paying the exponentially large size of
their background. In a unital $C^*$-algebra with positive normalized
trace $\tau$, let $H,Z$ be bounded selfadjoint elements. Duhamel's
ordered-simplex formula and tracial H\"older give
\begin{equation}\label{eq:duhamel-bound}
 \left|\frac{\mathrm d^r}{\mathrm dt^r}
          \tau(e^{H+tZ})\right|
 \le\norm Z^r\tau(e^{H+tZ}).
\end{equation}
The formula follows in the finite tracial von Neumann completion:
each product of exponential factors has H\"older exponents whose
reciprocals sum to one; the derivative factorial cancels the
simplex volume. The same argument for the first derivative gives
\[
 \tau(e^{H+tZ})\le e^{\abs t\norm Z}\tau(e^H).
\]
Here monotonicity is a statement about the traced exponential,
not operator monotonicity of the exponential map. Taylor's
remainder after degree $r-1$ is consequently bounded by
$\norm Z^r e^{\norm Z}\tau(e^H)/r!$.

Let $z$ be uniform on the circle, let $S_m$ be contractions, and
let $(b_m)$ be a finite real string indexed by positive integers. Put
\[
 X(z)=\sum_m b_m(z^mS_m+(\overline z)^mS_m^*),
 \qquad c=\sum_m\abs{b_m}\le4.
\]
Let $Y$ retain the first harmonic, or be zero. If
$d=\sum_{m\text{ omitted}}b_m^2$, the linear expectations vanish,
and unequal harmonics are orthogonal in the quadratic term.
For the remaining terms set
\[
 B_H(S)=\int_0^1\tau(e^{(1-v)H}S e^{vH}S^*)\dd v.
\]
Positivity and H\"older give $0\le B_H(S)\le\tau(e^H)$.
The two opposite orientations cancel the Taylor factor $1/2$;
their integrals agree by cyclicity. Thus the difference of the
quadratic expectations is exactly
$s^2\sum_{m\text{ omitted}}b_m^2B_H(S_m)$.
Charging the two cubic remainders yields, for real $s$,
\begin{equation}\label{eq:local-relative}
 \left|\E\tau(e^{H+sX})-\E\tau(e^{H+sY})\right|
 \le\left(s^2d+\frac83\abs s^3e^{8\abs s}c^3\right)\tau(e^H).
\end{equation}
Both expectations are at least $\tau(e^H)$ by convexity and
centering. Therefore this is a relative estimate in both directions,
uniform in the background $H$.

\subsection*{Every prime power, with its correct phase}

Give each prime an independent circle variable $z_p$, and replace
its local string by
\[
 X_p(z_p)=\sum_{m\log p<\ell}b_f(p^m)
       (z_p^mS_{m\log p}+(\overline{z_p})^mS_{-m\log p}).
\]
Unique factorization says that zero total displacement is equivalent
to zero signed exponent at each prime. The trace of every other
word is zero. The norm-convergent exponential series at fixed
$\ell$ therefore gives the exact identity
\begin{equation}\label{eq:phase-identity}
 \E\tau_\ell\left(e^{s\sum_pX_p(z_p)}\right)
       =\tau_\ell(e^{sW_{f,\ell}}).
\end{equation}
All harmonics at $p$ use the same $z_p$. Higher-order resonances
among these harmonics are retained, not assumed to vanish.

For good $p$ retain its first harmonic; for bad $p$ retain zero.
The local budgets satisfy
\begin{align}
 c_p&\le\frac{2\log p}{\sqrt p-1}\le4,
 &\sum_pc_p^3&\le C_*:=8\sum_p\frac{(\log p)^3}{(\sqrt p-1)^3}<\infty,
 \label{eq:c-budget}\\
 \sum_pd_p&\le D_*:=4\sum_p\frac{(\log p)^2}{p(p-1)}<\infty.
 \label{eq:d-budget}
\end{align}
At a bad prime its full squared mass is at most
$(\log p)^2/(p^2-1)$, so~\eqref{eq:d-budget} holds uniformly
even as the ramified set varies. Conditioning on all other primes
in~\eqref{eq:local-relative} and replacing one local string at a
time proves
\begin{equation}\label{eq:full-local-bound}
 \left|P_{f,\ell}(s)-\log\tau_\ell(e^{sW^{(1)}})\right|
 \le s^2D_*+\frac83\abs s^3e^{8\abs s}C_*.
\end{equation}
This is bounded independently of $f$ and $\ell$ at bounded $s$,
and tends to zero at $s=t/\ell$ for fixed $t$.
It is not a bound for $\norm{W-W^{(1)}}$.

\section{The weighted autocorrelation identity}\label{sec:profile}

For $u\in L^2(0,1)$ with $\norm u_2=1$ put
\[
 F(x)=\abs{u(x)}^2,\qquad
 c_u(a)=\int_0^{1-a}\overline{u(x)}u(x+a)\dd x,
 \qquad J(u)=\int_0^1a\abs{c_u(a)}^2\dd a.
\]
Let $\widehat F(n)=\int_0^1F(x)e^{-2\pi\ii nx}\dd x$.
We will use the exact identity
\begin{align}\label{eq:continuous-deficit}
 \frac1{12}-J(u)
 &=\frac1{2\pi^2}\sum_{n\ge1}\frac{\abs{\widehat F(n)}^2}{n^2}
 \notag\\
 &\quad+\frac12\int_0^1a\int_0^{1-a}\int_0^{1-a}
  \left|\overline{u(x)}u(x+a)-\overline{u(y)}u(y+a)\right|^2
             \dd x\dd y\dd a.
\end{align}
Both terms on the right are nonnegative.

To derive it, write $g_a(x)=\overline{u(x)}u(x+a)$. The elementary
variance identity on an interval of length $1-a$ is
\[
 (1-a)\int\abs{g_a}^2-\left|\int g_a\right|^2
   =\frac12\iint\abs{g_a(x)-g_a(y)}^2\dd x\dd y.
\]
The required integrations are legitimate without assuming $u\in L^4$.
Indeed Tonelli gives
\[
 \int_0^1a(1-a)\int_0^{1-a}F(x)F(x+a)\dd x\dd a
 =\frac12\int_0^1\int_0^1
   \abs{x-y}(1-\abs{x-y})F(x)F(y)\dd x\dd y<\infty.
\]
The periodic Bernoulli identity
\[
 v(1-v)=\frac16-\frac1{\pi^2}
                  \sum_{n\ge1}\frac{\cos(2\pi nv)}{n^2},
 \qquad 0\le v\le1,
\]
converges absolutely and uniformly. Integrating it against
$F(x)F(y)$ evaluates the preceding expression as
$1/12-(2\pi^2)^{-1}\sum_{n\ge1}\abs{\widehat F(n)}^2/n^2$.
Subtracting $J(u)$ proves~\eqref{eq:continuous-deficit}.

\subsection*{Equality and what the deficit actually controls}

Equality first forces $\widehat F(n)=0$ for $n\ne0$, hence
$F=1$ almost everywhere by Fourier uniqueness in $L^1$.
The second term then gives, for almost every $a$,
\[
 \overline{u(x)}u(x+a)=\chi(a)\quad\text{for almost every }x\in(0,1-a),
 \qquad\abs{\chi(a)}=1.
\]
The zero-extended $L^2$ autocorrelation is continuous, so
$\chi(a)=c_u(a)/(1-a)$ is continuous on $[0,1)$.
Continuity of translations in $L^2$ extends the increment identity
to every admissible $a$ on compact interior intervals.
Composing two increments gives
$\chi(a+b)=\chi(a)\chi(b)$ for $a,b\ge0$, $a+b<1$.
A continuous local circle character is $e^{\ii\omega a}$:
choose a continuous argument near zero, apply the additive law
there, and subdivide an arbitrary $a$ into small equal increments.
It follows that $e^{-\ii\omega x}u(x)$ has zero distributional
derivative in $(0,1)$ and is constant almost everywhere. Thus
\begin{equation}\label{eq:profile-equality}
 J(u)=\frac1{12}\quad\Longleftrightarrow\quad
 u(x)=e^{\ii\theta+\ii\omega x}\ \text{a.e.},\qquad \theta,\omega\in\R.
\end{equation}
If $J(u)\ge1/12-\varepsilon$, the first deficit is at most
$\varepsilon$, giving
$\sum_{n\ge1}\abs{\widehat F(n)}^2/n^2\le2\pi^2\varepsilon$.
This directly controls a negative-Sobolev measure of the squared
modulus. A stronger conclusion requires the following additional
compactness and local-coercivity arguments.

\subsection*{Strong stability modulo affine phase}

For arbitrary $u\in L^2(0,1)$ write
$D(u)=\norm u_2^4/12-J(u)\ge0$, and for unit $u$ let
\[
 \mathcal E=\{e^{\ii\theta+\ii\omega x}:\theta,\omega\in\R\},
 \qquad d(u)=\inf_{e\in\mathcal E}\norm{u-e}_2.
\]
The functional $J$ is a continuous real quartic polynomial on
the complex Hilbert space viewed as a real space. This follows
from $\abs{c_u(a)}\le\norm u_2^2$ and polarization.
Modulation and constant phase preserve $D$ and $d$.

Compactness modulo symmetries followed by transverse second
variation is a standard strategy in stability theory for sharp
Fourier inequalities; see, for example, Christ~\cite{Christ}.
Here a compact correlation splitting and explicit primitive
identities give a direct proof for this weighted functional.

We first show that $D(u_n)\to0$ for unit $u_n$ forces
$d(u_n)\to0$. The essential weak-splitting fact is
\begin{equation}\label{eq:weak-splitting}
 u_n\rightharpoonup u,\quad r_n=u_n-u
 \quad\Longrightarrow\quad J(u_n)=J(u)+J(r_n)+o(1).
\end{equation}
In the correlation expansion, each mixed term is a compact
linear or antilinear map of $r_n$ from $L^2_x$ to $L^2_a$.
For example its kernel is
$\overline{u(y-a)}\one_{\{0<a<y<1\}}$, of squared
Hilbert--Schmidt norm
$\int_0^1(1-x)\abs{u(x)}^2\dd x\le\norm u_2^2$.
These terms tend strongly to zero. The remaining cross term is
\[
 \int_{0<x<y<1}(y-x)c_u(y-x)
                         r_n(x)\overline{r_n(y)}\dd x\dd y.
\]
Its kernel is bounded, hence Hilbert--Schmidt, so it also tends
to zero. This proves~\eqref{eq:weak-splitting}; it does not
assume the self-correlation of a weak-null sequence tends to zero.

Plancherel for the zero-extended profile gives
\[
 J(u_n)\le\frac12\int_\R\abs{c_{u_n}(a)}^2\dd a
   =\frac1{4\pi}\int_\R\abs{\widehat{u_n}(\xi)}^4\dd\xi
   \le\frac12\norm{\widehat{u_n}}_\infty^2.
\]
Thus a near-maximizing sequence can be modulated and rotated so
that its mean is real and bounded below by a positive constant.
A weakly convergent subsequence then has a nonzero limit $u$.
If $\alpha=\norm u_2^2$, weak orthogonality gives
$\norm{r_n}_2^2\to1-\alpha$. By~\eqref{eq:weak-splitting}
and the homogeneous sharp inequality,
\[
 \frac1{12}\le\frac{\alpha^2+(1-\alpha)^2}{12}.
\]
Since $0<\alpha\le1$, necessarily $\alpha=1$. The convergence
is strong, and its limit belongs to $\mathcal E$ by
\eqref{eq:profile-equality}. This proves the claimed compactness
modulo modulation.

The local estimate has an explicit sharp quadratic constant.
Write $h=a+\ii b$ with $a,b$ real, $\mu=\int_0^1a$, and put
\[
 L_h(r)=\int_0^{1-r}\bigl(\overline{h(x)}+h(x+r)\bigr)\dd x.
\]
Exact expansion at the flat profile gives
\begin{equation}\label{eq:stability-expansion}
 D(1+h)=Q(h)+C_3(h)+D(h),\qquad
 \abs{C_3(h)}\le\frac73\norm h_2^3,
\end{equation}
where
\begin{align}\label{eq:Q-explicit}
 Q(h)&=\frac{\norm h_2^2}{6}+\frac{\mu^2}{3}
    -\int_0^1r\abs{L_h(r)}^2\dd r
    -2\int_0^1r(1-r)\Re c_h(r)\dd r,\notag\\
 C_3(h)&=\frac{\mu\norm h_2^2}{3}
    -2\Re\int_0^1rL_h(r)\overline{c_h(r)}\dd r.
\end{align}
The remainder bound follows from
$\abs{L_h(r)}\le2\norm h_2$ and
$\abs{c_h(r)}\le\norm h_2^2$.
Since $D\ge0$, $Q$ is nonnegative. Its zero directions contain
$\mathcal K=\operatorname{span}_{\R}\{1,\ii,\ii x\}$ by
scaling, phase, and modulation invariance. Positivity makes all
bilinear cross terms with $\mathcal K$ vanish.

In fact the transverse quadratic gap is exactly $1/6$.
If $h=a+\ii b\perp\mathcal K$ in the real Hilbert inner
product, then $\int a=\int b=\int xb=0$.
Let $A(x)=\int_0^xa$ and $B(x)=\int_0^xb$.
Use subscripts ${\rm e},{\rm o}$ for even and odd parts under
reflection $x\mapsto1-x$, and write $\overline A=\int_0^1A$.
Both primitives vanish at the endpoints, and
\[
 \Re L_h(r)=A(1-r)-A(r)=-2A_{\rm o}(r),\qquad
 \Im L_h(r)=-B(r)-B(1-r)=-2B_{\rm e}(r).
\]
For a real mean-zero function $v$ with primitive $V$, the same
Bernoulli expansion used above gives
\[
 -2\int_0^1r(1-r)c_v(r)\dd r
       =2\left(\norm V_2^2-\left|\int_0^1V\right|^2\right).
\]
This also follows from Parseval, since the nonzero Fourier
coefficients of $V$ are those of $v$ divided by $2\pi\ii n$.
The square of either reflection-parity component is symmetric,
so its integral against $r$ is half its integral. Substitution
in~\eqref{eq:Q-explicit}, using $\int B=-\int xb=0$, yields
\begin{equation}\label{eq:hessian-gap}
 Q(a)=\frac{\norm a_2^2}{6}
                   +2\norm{A_{\rm e}-\overline A}_2^2,
 \qquad
 Q(\ii b)=\frac{\norm b_2^2}{6}+2\norm{B_{\rm o}}_2^2.
\end{equation}
There are no real-imaginary cross terms. Hence
$Q(h)\ge\norm h_2^2/6$ on $\mathcal K^\perp$, with equality
for $h(x)=\cos(2\pi x)$.

For a unit $u$, a nearest affine phase exists because
\[
 d(u)^2=2-2\max_{\omega\in\R}
                 \left|\int_0^1u(x)e^{-\ii\omega x}\dd x\right|,
\]
and the Fourier transform is continuous, nonzero, and vanishes
at infinity. Align that nearest phase with $1$ and set $h=u-1$,
$\delta=\norm h_2=d(u)$. Differentiating the maximizing phase
and frequency gives $\Im\int h=\Im\int xh=0$; the derivatives
exist since $xu\in L^1$. Normalization gives
$\Re\int h=-\delta^2/2$. Thus
$\norm{P_{\mathcal K}h}_2^2=\delta^4/4$, and
\eqref{eq:stability-expansion}--\eqref{eq:hessian-gap} imply
\begin{equation}\label{eq:local-stability}
 D(u)\ge\delta^2\left(\frac16-\frac73\delta
                                        -\frac{\delta^2}{24}\right).
\end{equation}
For $\delta\le1/32$ this is at least $\delta^2/12$.
By the compactness already proved, the infimum of $D(u)$ over
unit profiles with $d(u)\ge1/32$ is some $\eta_0>0$.
Since $d(u)^2\le2$, the remaining profiles satisfy
$d(u)^2\le2D(u)/\eta_0$. This proves the global bound
in~\eqref{eq:stability-intro}, with an absolute but nonexplicit
constant $\max(12,2/\eta_0)$.

As $D(u)\downarrow0$, compactness places $u$ in the local
neighborhood, where $\delta^2\le12D(u)$.
Equation~\eqref{eq:local-stability} then gives uniformly
\begin{equation}\label{eq:sharp-stability}
 d(u)^2\le6D(u)+O(D(u)^{3/2}).
\end{equation}
This coefficient cannot be reduced. For small real $t$ take
$u_t=(1+t\cos(2\pi x))/\sqrt{1+t^2/2}$.
Positivity and the triangle inequality make frequency zero
maximize its Fourier modulus, so its nearest phase is $1$.
Then $d(u_t)^2\sim t^2/2$, whereas
\eqref{eq:stability-expansion} and~\eqref{eq:hessian-gap} give
$D(u_t)\sim t^2/12$. Their ratio tends to $6$. In particular
the square-root distance exponent is optimal. We make no
claim that $6$ is the sharp global constant.

\subsection*{The finite-dimensional identity}

For a unit vector $v=(v_0,\ldots,v_{d-1})$ define
$F_i=\abs{v_i}^2$,
$c_j=\sum_{i=0}^{d-j-1}\overline{v_i}v_{i+j}$, and
$\widehat F_d(k)=\sum_{i=0}^{d-1}F_i e^{-2\pi\ii ki/d}$.
The cyclic kernel $K_d(r)=r(d-r)$ has normalized Fourier coefficients
\[
 \widehat K_d(0)=\frac{d^2-1}{6},\qquad
 \widehat K_d(k)=-\frac1{2\sin^2(\pi k/d)},\quad 1\le k<d.
\]
For example these follow by applying the cyclic second difference:
$K_d(r+1)-2K_d(r)+K_d(r-1)=-2+2d\one_{\{r=0\}}$.
The same variance calculation on each overlap yields
\begin{align}\label{eq:discrete-deficit}
 \frac{d^2-1}{12}-\sum_{j=1}^{d-1}j\abs{c_j}^2
 &=\frac14\sum_{k=1}^{d-1}
       \frac{\abs{\widehat F_d(k)}^2}{\sin^2(\pi k/d)}\notag\\
 &\quad+\frac12\sum_{j=1}^{d-1}j
       \sum_{i,h=0}^{d-j-1}
       \abs{\overline{v_i}v_{i+j}-\overline{v_h}v_{h+j}}^2.
\end{align}
Only the kernel on the moduli has been evaluated cyclically;
$c_j$ is still the aperiodic, killed autocorrelation. In particular
\begin{equation}\label{eq:discrete-bound}
 \sum_{j=1}^{d-1}j\abs{c_j}^2\le\frac{d^2-1}{12}.
\end{equation}
For $d\ge2$, equality forces $F_i=1/d$, and equality at lag one
forces all successive phase increments to agree. Its complete equality class
is $v_i=d^{-1/2}e^{\ii\theta+\ii\omega i}$. The $d=1$ case is
immediate.

Equivalently, every nonnegative trigonometric polynomial $Q$ of
degree at most $n$ and mean one satisfies
\[
 \sum_{k\in\mathbb Z}\abs k\abs{\widehat Q(k)}^2
           \le\frac{n(n+2)}6.
\]
Fej\'er--Riesz factorization gives this formulation from
\eqref{eq:discrete-bound}; its extremizers for $n\ge1$ are translated
Fej\'er kernels. This is a useful equivalent description of the
same inequality, not an additional independent extremal result.

For comparison with a weighted linear positive-definite extremum,
$\Lambda(a)=c_u(a)$ is positive definite and supported in $[-1,1]$,
and $\abs\Lambda^2$ is positive definite as well. Enlarging the
class of $\abs\Lambda^2$ to all normalized positive-definite
functions loses the square constraint. That larger class contains
$\phi(a)=(1-\abs a)_+$, for which
$\int_{-1}^1\abs a\phi(a)\dd a=1/3$; our quadratic optimum is
$\int_{-1}^1\abs a(1-\abs a)^2\dd a=1/6$.
Thus the direct application of a linear weighted theorem such
as~\cite{Manov} cannot give the constant in~\eqref{eq:variational-intro}.

\section{Uniform pressure through a boundary-preserving grid}\label{sec:pressure}

We prove~\eqref{eq:uniform-main}. Fix $C>0$ and take
$2\le t\le C\ell$; negative $t$ will follow at the end.
All sums in this section initially contain only good first primes.

\subsection*{The flat lower bound}

After prime-phase expectation the exponential operator is a
multiplication operator: exactly the primewise balanced words
survive, and every such word is multiplication by its full
survival indicator. Its coefficients are nonnegative products of
even powers of the real $b_p$. Taking the flat unit vector
$v=\ell^{-1/2}\one_{I_\ell}$ and applying scalar Jensen to its
spectral measure gives
\begin{equation}\label{eq:bessel-lower}
 \tau_\ell(e^{tW^{(1)}/\ell})
 \ge\prod_{h_p<\ell}I_0\left(
          \frac{2t b_p}{\ell}(1-h_p/\ell)\right),
\end{equation}
where $I_0(2x)=\sum_{r\ge0}x^{2r}/(r!)^2$.
For bounded $\abs x$, $\log I_0(2x)=x^2+O(x^4)$.
Using~\eqref{eq:variance} and integration by parts,
\[
 \sum_{h_p<\ell}b_p^2(1-h_p/\ell)^2
   =\frac2\ell\int_0^\ell V_f(H)(1-H/\ell)\dd H
   =\frac{\ell^2}{12}+O_f(1).
\]
The fourth-power remainder is $O_C(t^4/\ell^4)=O_C(1)$.
This proves the lower side of~\eqref{eq:uniform-main} for
$W^{(1)}$, with a bounded error throughout the stated range.

\subsection*{Word truncation and rounding}

Choose once and for all $D>0$ with $D\log2>4$, and set
\begin{equation}\label{eq:mesh}
 M=\lceil Dt^2\rceil,\quad q=\lceil t^6\rceil,\quad
 d=q+2M,\quad j_p=\lfloor qh_p/\ell\rfloor.
\end{equation}
Here $0\le j_p<q$; the zero bin must be retained.
On $\C^d$ use the killed integer shifts, with $S_0=I_d$, and
give each prime its own independent circle phase, even when
several primes have the same $j_p$. Put
\[
 B_z=\sum_{h_p<\ell}\frac{b_p}{\ell}
                 (z_pS_{j_p}+\overline z_pS_{j_p}^*).
\]
A balanced word of length $r\le M$ has original rescaled range
$R$ and integer-grid range $R_q$. Each partial sum moves by at
most $r/q$, hence $\abs{R_q/q-R}\le2r/q$. It is exactly closed
on both grids, and
\begin{equation}\label{eq:mask-domination}
 (1-R)_+\le\frac dq(1-R_q/d)_+.
\end{equation}
Indeed $d-R_q\ge q(1-R)+2(M-r)$. Unbalanced arithmetic words
which close after rounding still vanish under the independent
prime phases.

The positive first-prime series has tail above degree $M$ at most
\[
 2^{-M}\prod_p I_0(4tb_p/\ell)
 \le\exp\{-M\log2+4t^2V_f(\ell)/\ell^2\}
 \le e^{-t^2}
\]
for all sufficiently large $\ell$, uniformly in $2\le t\le C\ell$.
We used $I_0(2x)\le e^{x^2}$ and
$V_f(\ell)/\ell^2\to1/2$. Consequently
\begin{equation}\label{eq:grid-comparison}
 \tau_\ell(e^{tW^{(1)}/\ell})
       \le\frac dq\E_z\tr_d(e^{tB_z})+e^{-t^2},
 \qquad \tr_d=\Tr/d.
\end{equation}
This multiplicative comparison is uniform at the exponential
coupling being estimated; it is not a fixed-moment approximation.

\subsection*{Relative replacement by circular Gaussians}

Replace $z_p$ successively by independent
$\zeta_p=(g_{p,1}+\ii g_{p,2})/\sqrt2$, where the $g_{p,a}$ are
real standard Gaussians. Denote the resulting matrix by $G$.
Circle and circular Gaussian variables have the same mixed
moments through degree three. For a Hermitian background $H$,
the fourth-order version of~\eqref{eq:duhamel-bound} bounds the
Taylor remainder by
$\norm Z^4e^{\norm Z}\Tr(e^H)/24$.
With $Z=(tb_p/\ell)(zS+\overline zS^*)$, its expected value is
at most $C_C(t b_p/\ell)^4\Tr(e^H)$, for either distribution.
This follows from $\norm Z\le2\abs{tb_p/\ell}\abs z$ and
the finite Gaussian exponential moments; $\abs{tb_p/\ell}$
is bounded uniformly by a constant depending only on $C$.
Both centered expectations are at least $\Tr(e^H)$. As in the
local restoration argument, comparison is relative and gives
\begin{equation}\label{eq:gauss-replacement}
 \left|\log\E_z\tr_d(e^{tB_z})-
             \log\E\tr_d(e^{tG})\right|
 \le C_C\frac{t^4}{\ell^4}\sum_pb_p^4=O_C(1).
\end{equation}
There is no dimension-dependent denominator or norm bound on the
background in this estimate.

\subsection*{The two Gaussian variances}

Write $G=\sum_a g_aA_a$. Direct multiplication gives
\[
 \sum_aA_a^2
  =\ell^{-2}\sum_pb_p^2(S_{j_p}S_{j_p}^*+S_{j_p}^*S_{j_p})
  \le2V_f(\ell)\ell^{-2}I_d.
\]
The Gaussian matrix-series bound~\cite{Tropp} therefore yields
\begin{equation}\label{eq:gaussian-mean}
 \E\lambda_{\max}(G)
       \le2\sqrt{V_f(\ell)\ell^{-2}\log d}
       =O_f(\sqrt{\log(2+t)}).
\end{equation}
For completeness, its trace-mgf estimate follows by Gaussian
integration by parts and
\[
 \Tr(Ae^{vB}Ae^{(1-v)B})\le\Tr(A^2e^B),\quad0\le v\le1,
\]
for Hermitian $A,B$, which is scalar convexity in an eigenbasis
of $B$. Thus $\E\Tr e^{\theta G}\le
d\exp(\theta^2\norm{\sum A_a^2}/2)$; optimizing in $\theta$
proves~\eqref{eq:gaussian-mean}.

The sharp quadratic coefficient uses a different quantity:
\begin{equation}\label{eq:weak-variance}
 \sigma_w^2=\sup_{\norm v=1}\sum_a\ip v{A_av}^2
     =\frac2{\ell^2}\sup_{\norm v=1}
                    \sum_pb_p^2\abs{c_{j_p}(v)}^2,
 \qquad c_0(v)=1.
\end{equation}
The top eigenvalue is the supremum of its real linear Rayleigh
functionals, and so is $\sigma_w$-Lipschitz in the real Gaussian
variables. Gaussian logarithmic Sobolev concentration~\cite{Gross,Ledoux}
gives
\begin{equation}\label{eq:weak-concentration}
 \log\E e^{t\lambda_{\max}(G)}
       \le t\E\lambda_{\max}(G)+t^2\sigma_w^2/2.
\end{equation}
In detail, for a Lipschitz function $F$ with constant $\sigma_w$,
the logarithmic Sobolev inequality applied to $e^{tF/2}$ gives
$t\Psi'(t)-\Psi(t)\le t^2\sigma_w^2/2$, where
$\Psi(t)=\log\E e^{tF}$; integrate from zero. Smooth
approximation proves the assertion for the top eigenvalue.

\subsection*{All grid vectors and all intermediate couplings}

Let $H_0=B_f\log^2(2+\ell)$, where $B_f$ is large enough that
$c_f\sqrt{B_f}>12$, and take $\ell$ so large that $H_0<\ell$.
Write $\Delta=\ell/q$. All bins whose left endpoint is less
than $H_0$ have total normalized mass at most
\begin{equation}\label{eq:short-bin-budget}
 \ell^{-2}V_f(H_0+\Delta)
       =O_f\left(\frac{\log^4(2+\ell)}{\ell^2}+q^{-2}\right).
\end{equation}
This includes the scalar zero bin and any straddling bin.
For the remaining bins,~\eqref{eq:bin} and $q=O_C(\ell^6)$
show that the sum of all normalized errors is $O_{f,C}(\ell^{-6})$.
Their main masses are $(j+1/2)/q^2$.
Since every $\abs{c_j(v)}\le1$, the error bounds hold for all
unit vectors simultaneously. Enlarge the positive main sum to
$1\le j<d$ and apply~\eqref{eq:discrete-bound}. We obtain
\begin{align}\label{eq:uniform-weak-bound}
 \sigma_w^2
 &\le\frac2{q^2}\left(\frac{d^2-1}{12}+\frac d2\right)
      +O_{f,C}\left(\frac{\log^4(2+\ell)}{\ell^2}+q^{-2}
                          +\ell^{-6}\right)\notag\\
 &=\frac16+O_{f,C}\left(\frac Mq+\frac1q+
                \left(\frac Mq\right)^2+
                \frac{\log^4(2+\ell)}{\ell^2}\right).
\end{align}
The extra $q^{-2}$ and $\ell^{-6}$ terms are absorbed in the
last line. This is where weak convergence of the variance measures
would have been insufficient: the maximizing vector changes with
the grid, and all its bins must be controlled together.

Since $\tr_d(e^{tG})\le e^{t\lambda_{\max}(G)}$, equations
\eqref{eq:gaussian-mean}--\eqref{eq:uniform-weak-bound} give
\[
 \log\E\tr_d(e^{tG})
 \le\frac{t^2}{12}+O_{f,C}\left(
       t\sqrt{\log(2+t)}+1+
                \frac{t^2\log^4(2+\ell)}{\ell^2}\right).
\]
Here $t^2(M/q+q^{-1}+(M/q)^2)=O_D(1)$.
The last displayed error is uniformly absorbed by
$O_C(t\sqrt{\log(2+t)}+1)$ when $2\le t\le C\ell$:
indeed $t\log^4(2+\ell)/\ell^2\le
C\log^4(2+\ell)/\ell$, which is bounded.
The factor $d/q$, the additive tail in~\eqref{eq:grid-comparison},
and~\eqref{eq:gauss-replacement} cost only bounded amounts in
the logarithm. The Gaussian trace expectation is at least one,
so the additive tail creates no denominator difficulty.

Finally~\eqref{eq:full-local-bound} restores the complete local
strings at coupling $t/\ell\le C$, again at bounded cost.
The first-prime trace is even in $t$, since replacing every
$z_p$ by $-z_p$ changes its sign. Apply the same restoration for
negative $t$. Together with~\eqref{eq:bessel-lower}, this proves
\eqref{eq:uniform-main} in its full range, and hence
\eqref{eq:pressure-main} and~\eqref{eq:K-main}.

\section{Intermediate scales and boundary survival}\label{sec:deviations}

\subsection*{A scalar large-deviation argument}

Let $r_\ell\to\infty$ with $r_\ell=O(\ell)$, and let $\mu_\ell$
be the probability spectral law of $W_{f,\ell}/(\ell r_\ell)$
in $\tau_\ell$. Equation~\eqref{eq:uniform-main} gives, for every
fixed $s\in\R$,
\begin{equation}\label{eq:scaled-mgf}
 \frac1{r_\ell^2}\log\int e^{sr_\ell^2x}\dd\mu_\ell(x)
   =\frac1{r_\ell^2}P_{f,\ell}(sr_\ell/\ell)
   \longrightarrow\frac{s^2}{12}.
\end{equation}
For $s\ne0$ one applies the uniform estimate at $t=\abs s r_\ell$
with its fixed upper constant enlarged by $\abs s$; $s=0$ is
immediate. This works however slowly $r_\ell$ diverges.

We give the short scalar tilting proof to specify the measure
and the quantifiers. The same argument will later apply to
the central limiting law. More generally let probability measures
$\mu_j$ have speed $v_j\to\infty$ and scaled logarithmic mgf
$L_j(s)=\log\int e^{sv_jx}\dd\mu_j(x)$ with
$L_j(s)/v_j\to s^2/12$ for every real $s$. Define
\[
 \dd\mu_{j,s}(x)=e^{sv_jx-L_j(s)}\dd\mu_j(x).
\]
For $h>0$, exponential Markov gives
\[
 \limsup_j v_j^{-1}\log\mu_{j,s}(x\ge s/6+\eta)
 \le\frac{(s+h)^2-s^2}{12}-h(s/6+\eta)
 =h^2/12-h\eta.
\]
Take $h=6\eta$. A negative increment proves the corresponding
left-tail estimate. Thus
\begin{equation}\label{eq:tilted-concentration}
 \limsup_jv_j^{-1}\log\mu_{j,s}(\abs{x-s/6}\ge\eta)
       \le-3\eta^2.
\end{equation}
For any $a\in\R$, set $s=6a$ and undo the tilt on
$(a-\eta,a+\eta)$; its tilted mass tends to one. This gives
\[
 \liminf_jv_j^{-1}\log\mu_j((a-\eta,a+\eta))
       \ge-3a^2-6\abs a\eta.
\]
Shrinking the interval inside an open set proves the lower bound.
For an upper bound near $a$ use Markov with $s=6a$; a finite
cover gives the bound on compact sets. Markov at $s=\pm6R$
also gives exponential tightness,
$\limsup_jv_j^{-1}\log\mu_j(\abs x\ge R)\le-3R^2$.
It extends the compact upper bound to every closed set.
Consequently for every Borel set $B$,
\begin{equation}\label{eq:LDP}
 -\inf_{B^\circ}3x^2
 \le\liminf_jv_j^{-1}\log\mu_j(B)
 \le\limsup_jv_j^{-1}\log\mu_j(B)
 \le-\inf_{\overline B}3x^2.
\end{equation}
We use $\log0=-\infty$ and $\inf\varnothing=+\infty$.
Applying this proof to~\eqref{eq:scaled-mgf} establishes the
claimed intermediate-scale principle. For every fixed $a>0$,
\begin{equation}\label{eq:moderate-tails}
 \begin{split}
 \lim_{\ell\to\infty}r_\ell^{-2}
   \log\tau_\ell(\one_{[a\ell r_\ell,\infty)}(W_{f,\ell}))&=-3a^2,\\
 \lim_{\ell\to\infty}r_\ell^{-2}
   \log\tau_\ell(\one_{(-\infty,-a\ell r_\ell]}(W_{f,\ell}))&=-3a^2.
 \end{split}
\end{equation}
The equality follows by sandwiching between open and closed rays;
their rate infima agree. Equal leading tail rates do not assert
exact symmetry of the full finite-window prime-power law.

\subsection*{Which spectral mass carries the Weyl coefficient}

Taking $r_\ell=\ell$, the exponential tilt at fixed coupling $s$
concentrates $W_{f,\ell}/\ell^2$ at $s/6$ by
\eqref{eq:tilted-concentration}. At the Weyl coupling $s=1/2$
this value is $1/12$.
One can also recover the tilted first moment without differentiating
an asymptotic expansion. The functions
$\ell^{-2}P_{f,\ell}(s)$ are convex and differentiable. For $h>0$
their derivatives lie between their left and right difference
quotients. First let $\ell$ tend to infinity in these bounds and
then let $h\downarrow0$. This proves
\begin{equation}\label{eq:tilted-mean}
 \frac{\tau_\ell(W_{f,\ell}e^{sW_{f,\ell}})}
        {\ell^2\tau_\ell(e^{sW_{f,\ell}})}\longrightarrow\frac s6.
\end{equation}
No convergence of second derivatives, tilted variances, or
tilted Gaussian fluctuations is inferred by this argument.

\subsection*{The exact exponential cost of surviving the boundary}

For first primes and fixed $s\ne0$, introduce independent integer
counts $R_p$ with
\[
 \mathbb P(R_p=r)=\frac{(\abs s\abs{b_p})^{2r}}
                         {(r!)^2I_0(2\abs s\abs{b_p})},\qquad r\ge0.
\]
Given these counts, choose uniformly an ordering of $R_p$ copies
of $+h_p$ and $R_p$ copies of $-h_p$ for each prime. Let
$F_\ell=(1-\operatorname{range}/\ell)_+$ be the surviving
fraction of starting positions. Counting these orderings in the
exponential series gives the exact formula
\begin{equation}\label{eq:survival-identity}
 \tau_\ell(e^{sW^{(1)}})
       =\prod_pI_0(2\abs s\abs{b_p})\,\E F_\ell.
\end{equation}
These are paired words, not paths with independent increments.
By~\eqref{eq:variance} and the summable fourth-power error,
the logarithm of the unmasked product is
$s^2\ell^2/2+O_{f,s}(1)$. The pressure estimate gives
\begin{equation}\label{eq:survival-rate}
 \log\E F_\ell=-\frac{5s^2}{12}\ell^2
                  +O_{f,s}(\ell\sqrt{\log(2+\ell)}).
\end{equation}
Thus the boundary contributes at leading exponential order,
not merely as a lower-order endpoint correction.

\section{The central arithmetic law and its sharp tails}\label{sec:central}

\subsection*{A complete moment description}

For a pair partition $\pi$ of $\{1,\ldots,2m\}$ label its
pairs in order of first occurrence. Give pair $j$ a length
$a_j\in[0,1]$ and an orientation $\varepsilon_j\in\{1,-1\}$.
Its first occurrence contributes increment $\varepsilon_ja_j$
and its second contributes $-\varepsilon_ja_j$. If
$r_0=0,r_1,\ldots,r_{2m}=0$ are the partial sums, put
\[
 F_{\pi,\varepsilon}(a)
    =(1-\max_i r_i+\min_i r_i)_+.
\]
The central moments are
\begin{equation}\label{eq:central-moments}
 m_0=1,\quad m_{2m+1}=0,\qquad
 m_{2m}=\sum_\pi\sum_{\varepsilon\in\{\pm1\}^m}
        \int_{[0,1]^m}F_{\pi,\varepsilon}(a)
                           \prod_{j=1}^m a_j\dd a_j.
\end{equation}
In particular $0\le m_{2m}\le(2m-1)!!$, since the masks are
at most one and $2\int_0^1a\dd a=1$. The formal exponential
series therefore satisfies
\begin{equation}\label{eq:central-entire}
 M(s):=\sum_{m\ge0}\frac{m_{2m}s^{2m}}{(2m)!}\le e^{s^2/2}
 \qquad(s\in\R).
\end{equation}

We obtain existence of its probability measure from the actual
operators. Expand $\tau_\ell((W^{(1)}/\ell)^r)$. Every surviving
word is balanced at each prime, so odd moments vanish. At degree
$2m$, words with exactly $m$ distinct primes, each occurring twice,
are indexed by $\pi$, $\varepsilon$, and distinct prime labels.
Their masks are exactly those in~\eqref{eq:central-moments}.
Write $c_{p,\ell}=b_p/\ell$. Allowing repeated labels in this
pairing formula, and bounding words with a prime used four or
more times, costs at most
\[
 C_m\left(\sum_pc_{p,\ell}^4\right)
           \left(1+\sum_pc_{p,\ell}^2\right)^{m-2}=o_f(1)
 \quad(m\ge2).
\]
To see this, select two paired labels with a common prime,
sum that prime with fourth-power weight, and sum the other
labels freely. There are only finitely many overcountings at
fixed $m$, and every mask is bounded by one. For $m=1$ there
is no collision error.

The masks are continuous in their lengths. The variance-measure
convergence in~\eqref{eq:variance-measure} now gives exactly
\eqref{eq:central-moments}. Also the unmasked positive-word bound
gives, for every real $s$,
\begin{equation}\label{eq:central-UI}
 \tau_\ell(e^{sW^{(1)}/\ell})
 \le\prod_pI_0(2sb_p/\ell)
 \le\exp(s^2V_f(\ell)/\ell^2).
\end{equation}
Second moments give tightness of the spectral laws. Higher moments
give uniform integrability of each fixed power, and
\eqref{eq:central-UI} at larger positive and negative arguments
gives uniform integrability of each fixed exponential.
Every subsequential limit has~\eqref{eq:central-moments}, and
the Gaussian moment majorant makes this moment problem determinate.
Hence there is a unique probability measure $\nu$, with mgf
$M$, and the whole sequence converges weakly, in moments, and
in real exponential moments.

Applying~\eqref{eq:full-local-bound} at $s/\ell$ shows that the
full and first-prime logarithmic mgfs differ by $O_s(\ell^{-2})$.
It also provides uniform exponential bounds at all larger fixed
arguments for the full operators. Tightness, subsequence
uniqueness, and uniform integrability then prove the same three
convergences for $W_{f,\ell}/\ell$. All local factors are included.

\subsection*{The Toeplitz identification and the low moments}

For the matrices~\eqref{eq:toeplitz-intro}, the circular Gaussian
Wick formula pairs opposite orientations. A closed integer word
has normalized trace $(1-\operatorname{range}/n)_+$.
The variance of its $j$th diagonal is $j/n^2$, and the sums over
$j/n$ are Riemann sums. The expected normalized trace moments
therefore converge to~\eqref{eq:central-moments}. Uniform
unmasked Gaussian bounds justify all fixed exponential moments
as well. This identifies $\nu$ as the expected limiting spectral
measure; no almost-sure assertion is needed.

This law is the $\sqrt a$ specialization of the variance-profile
theorem in~\cite{BSS}, Theorem 3.3 and Section 7.2.2 in the
arXiv version cited there.
Their real symmetric formulation has the same limiting moments.
Indeed a real Wick pair with equal rather than opposite
orientations imposes a nontrivial linear closure equation on
the diagonal lengths. At fixed order this leaves $O(n^{m-1})$
choices rather than $O(n^m)$; each variance is at most $1/n$,
so its normalized contribution is $O(n^{-1})$.
Coincident diagonal lengths have the same lower-order counting
loss. The opposite orientations give exactly the circular
complex limit above.

At order two the two orientations each have mask $1-a$, giving
$m_2=2\int_0^1a(1-a)\dd a=1/3$. At order four set
\[
 I_{\max}=\int_0^1\int_0^1ab(1-\max(a,b))\dd a\dd b=\frac1{20},
 \qquad
 I_{\Sigma}=\int_{a+b\le1}ab(1-a-b)\dd a\dd b=\frac1{120}.
\]
Splitting the square into its two triangles proves the first
evaluation; integrating in $b$ first proves the second.
For each of the pair words $AABB$ and $ABBA$, two orientations
give $I_{\max}$ and two give $I_{\Sigma}$. The word $ABAB$
has four orientations giving $I_{\Sigma}$. Therefore
\begin{equation}\label{eq:moments-low}
 m_2=\frac13,\qquad m_4=4I_{\max}+8I_{\Sigma}=\frac4{15}.
\end{equation}
These numbers are evaluations of the specified profile law,
not separate general random-matrix assertions.

\subsection*{A direct large-coupling estimate for the limiting law}

Taking $\ell\to\infty$ at fixed $s$ in the flat-vector lower
bound~\eqref{eq:bessel-lower} gives
\begin{equation}\label{eq:central-lower}
 M(s)\ge e^{s^2/12}\qquad(s\in\R).
\end{equation}
For a matching upper estimate work directly with
\eqref{eq:central-moments}, not with a diagonal substitution
in a fixed-$s$ convergence statement. By symmetry take $s\ge2$.
Choose $L=\lceil Ds^2\rceil$, $q=\lceil s^6\rceil$,
$d=q+2L$, with $D\log2>4$. The part of the positive moment
series above degree $L$ is at most
$2^{-L}M(2s)\le e^{-s^2}$ by~\eqref{eq:central-entire}.
Round each length $a$ down to $j/q$. The same enlarged-interval
argument as in~\eqref{eq:mask-domination} dominates every
retained mask by $d/q$ times its grid mask. The bin masses now
are exact:
\[
 w_j=\int_{j/q}^{(j+1)/q}a\dd a=\frac{j+1/2}{q^2},
 \qquad 0\le j<q.
\]
Let
\[
 G_{q,d}=\sum_{j=0}^{q-1}\sqrt{w_j}
                  (z_jS_j+\overline z_jS_j^*),\qquad S_0=I_d.
\]
The Wick expansion of its exponential is precisely the rounded
pairing expansion, including multiple pairs in one bin and the
scalar zero-lag term. Thus
\begin{equation}\label{eq:central-grid}
 M(s)\le\frac dq\E\tr_d(e^{sG_{q,d}})+e^{-s^2}.
\end{equation}
The strong real-Gaussian variance is at most
$2\sum_jw_jI_d=I_d$, so $\E\lambda_{\max}(G_{q,d})
\le\sqrt{2\log d}$. By~\eqref{eq:discrete-bound} its weak variance
is bounded by
\[
 \frac2{q^2}\left(\frac{d^2-1}{12}+\frac d2+\frac12\right)
       =\frac16+O(L/q+q^{-1}+(L/q)^2).
\]
The Gaussian concentration calculation now gives
\begin{equation}\label{eq:central-pressure}
 \frac{s^2}{12}\le\log M(s)
       \le\frac{s^2}{12}+C\abs s\sqrt{\log(2+\abs s)}
       \quad(\abs s\ge s_0),
\end{equation}
with absolute constants. The grid errors multiplied by $s^2$
are bounded, $d/q$ is bounded, and the additive tail in
\eqref{eq:central-grid} is harmless because the centered
Gaussian trace expectation is at least one.

If $Y$ has law $\nu$, then
$R^{-2}\log\E e^{sR^2(Y/R)}=R^{-2}\log M(sR)\to s^2/12$.
Apply the tilting proof of Section~\ref{sec:deviations} with
speed $R^2$. Open and closed rays give exactly the two limits
in~\eqref{eq:central-intro}. In particular both sides of the
support of $\nu$ are unbounded. The quadratic logarithmic
tails do not make the distribution Gaussian.

\section{Several arithmetic forms in one translation algebra}\label{sec:colored}

Fix pairwise inequivalent primitive holomorphic forms
$f_1,\ldots,f_d$, each satisfying the earlier real,
trivial-character hypotheses. The tuple is fixed throughout.
Let $\mathcal P_0$ be the primes good for every member, and put
$b_i(p)=\lambda_{f_i}(p)\log p/\sqrt p$ there.
The off-diagonal version of the same prime number theorem,
\cite[Theorem 1.1]{GJ}, gives
\begin{equation}\label{eq:cross-RS}
 \sum_{\substack{p\le X\\p\in\mathcal P_0}}
        \lambda_{f_i}(p)\lambda_{f_j}(p)\log p
 =\delta_{ij}X+O_{\boldsymbol f}(Xe^{-c_{\boldsymbol f}\sqrt{\log X}}).
\end{equation}
The theorem excludes imaginary norm twists in its off-diagonal
case. Here such an equivalence would force the twist parameter
to be zero because both central characters are trivial; thus
pairwise inequivalence suffices. Distinct quadratic twists qualify
only when the resulting representations are genuinely inequivalent.
A complex-multiplication self-twist does not create a new color.

Weighted partial summation yields
$\sum_{h_p<H}b_i(p)b_j(p)=\delta_{ij}H^2/2+O_{\boldsymbol f}(1)$.
The signed measures
$\ell^{-2}\sum b_i(p)b_j(p)\delta_{h_p/\ell}$ converge to
$\delta_{ij}a\dd a$ and have uniformly bounded total variations
by Cauchy--Schwarz. Their product measures consequently converge
against continuous masks: approximate a mask uniformly by sums
of product functions and use the total-variation bounds.

\subsection*{The colored moment formula}

Put $X_{i,\ell}=W_{f_i,\ell}/\ell$, all in the same algebra
$\mathcal A_\ell$. The limit of a fixed mixed moment of even
length is
\begin{equation}\label{eq:colored-law}
 \lim_{\ell\to\infty}\tau_\ell(X_{i_1,\ell}\cdots X_{i_{2m},\ell})
 =\sum_{\substack{\pi\text{ a pairing}\\
          \text{each pair has equal colors}}}
       \sum_{\varepsilon\in\{\pm1\}^m}
       \int F_{\pi,\varepsilon}(a)\prod_{j=1}^m a_j\dd a_j.
\end{equation}
Odd limits vanish. For the common-good first-prime parts the
proof is the earlier distinct-prime pairing expansion, now using
the signed cross-variance measures. Repeated labels are bounded
by the same fourth-power majorant, since the tuple is fixed.

To restore all local factors in mixed words, write $R_i$ for
the difference between the full operator and its common-good
first-prime part. The two-letter trace gives exactly
\[
 \tau_\ell(R_i^2)=2\sum_{\text{omitted active }p^m}
               b_{f_i}(p^m)^2(1-m\log p/\ell)_+
       \le C_{\boldsymbol f}.
\]
Besides higher powers this includes finitely many primes bad
for another tuple member. Thus
$\norm{R_i/\ell}_{L^2(\tau_\ell)}=O_{\boldsymbol f}(\ell^{-1})$.
The scalar exponential bounds give uniformly bounded tracial
$L^p$ norms at each fixed finite $p$ for all normalized full and
first-prime operators. Telescope a mixed word of length $r\ge2$,
put its remainder in $L^2$, and its other $r-1$ factors in
$L^{2(r-1)}$. Tracial H\"older makes each replacement
$O_{\boldsymbol f,r}(\ell^{-1})$. The first moments are zero.
This proves~\eqref{eq:colored-law} for the original full operators.

Taking limits of $\tau_\ell(P^*P)$ shows that these moments
define a positive normalized tracial functional on polynomials
in selfadjoint variables $X_1,\ldots,X_d$. It is the weighted
colored Toeplitz law realized by independent circular Gaussian
diagonal families on one common matrix space. This does not
assert the existence of bounded limiting variables; their scalar
laws have unbounded support.

For distinct colors $X,Y$, the pair integrals already computed give
\begin{equation}\label{eq:mixed-moments}
 \tau(XY)=0,\qquad
 \tau(X^2Y^2)=\frac7{60},\qquad
 \tau(XYXY)=\frac1{30},\qquad
 \tau([X,Y]^*[X,Y])=\frac16.
\end{equation}
The last identity follows by expanding the commutator and using
cyclicity. Thus orthogonality of the arithmetic prime coefficients
does not produce commuting variables or freeness. Nor are the
second powers classically independent: $7/60\ne1/9$.
These profile-specific values illustrate the established colored
Toeplitz phenomenon~\cite{BHS}.

\subsection*{Orthogonal invariance without classical independence}

For $a\in\R^d$ set $X(a)=\sum_i a_iX_i$. Multilinearity in
\eqref{eq:colored-law} gives, for vectors $a_1,\ldots,a_{2m}$,
\begin{equation}\label{eq:vector-pairing}
 \tau(X(a_1)\cdots X(a_{2m}))
 =\sum_\pi\left(\prod_{\{r,s\}\in\pi}a_r\cdot a_s\right)
        \sum_{\varepsilon}\int F_{\pi,\varepsilon}(v)
                                        \prod_jv_j\dd v_j.
\end{equation}
Hence the entire joint moment law is invariant under real
orthogonal changes of color. In particular the scalar law of
$X(a)$ is the dilation of $\nu$ by $\norm a_2$; its moments
and exponential moments determine it uniquely. Orthogonal
invariance here is compatible with the nonzero commutator
in~\eqref{eq:mixed-moments}; it is not a multivariate classical
Gaussian law.

There is also a growing-window statement for each fixed real
linear combination $B_{a,\ell}=\sum_i a_iW_{f_i,\ell}$:
\begin{equation}\label{eq:vector-pressure}
 \log\tau_\ell(e^{B_{a,\ell}})
   =\frac{\norm a_2^2}{12}\ell^2
                +O_{\boldsymbol f,a}(\ell\sqrt{\log(2+\ell)}),
\end{equation}
with bounded lower error. Retain every actual first harmonic
and form the real coefficient array $\sum_i a_i b_i(p)$.
Equation~\eqref{eq:cross-RS} gives its variance density
$\norm a_2^2 a\dd a$, including the moving-bin error.
The positive balanced-word and Gaussian arguments then apply
to this single array. The higher local harmonics are restored
by~\eqref{eq:local-relative}, with constants depending on the
fixed tuple and $a$. One phase is used at each prime, shared
by all the forms and all its harmonics. Assigning independent
phases to different forms at the same prime would change the
operator and is not done.

This is a family of scalar spectral assertions in a noncommutative
algebra. It does not define a joint classical spectral measure
for all $W_{f_i,\ell}$, and makes no estimate uniform in a
growing tuple or a growing conductor.

\section{Tapers with equal variance and different Weyl growth}\label{sec:tapers}

The constant $1/12$ reflects where the variance lies among the
translation lengths, not just its total size. We make this
distinction quantitative for an explicit deformation of the
actual arithmetic operator.

\subsection*{A boundary-profile criterion}

Let $w\ge0$ be continuous on $[0,1]$, and set
\[
 J_w(u)=\int_0^1w(a)\abs{c_u(a)}^2\dd a,\qquad
 K(a)=w(a)(1-a).
\]
Assume $K(0)=K(1)=0$, $K(a)=K(1-a)$, and that the periodic
Fourier series of $K$ is absolutely convergent with
$\widehat K(n)\le0$ for every $n\ne0$.
The coefficients are real by symmetry. The overlap
Cauchy--Schwarz argument and Fourier expansion give, for unit $u$,
\begin{align}\label{eq:general-profile}
 J_w(u)&\le\frac12\int_0^1\int_0^1
                 K(x-y)\abs{u(x)}^2\abs{u(y)}^2\dd x\dd y\notag\\
 &=\frac{\widehat K(0)}2+
              \sum_{n\ge1}\widehat K(n)\abs{\widehat F(n)}^2
 \le\frac{\widehat K(0)}2.
\end{align}
The symmetry of $K$ implies
\begin{equation}\label{eq:profile-constant}
 c(w):=\frac{\widehat K(0)}2
       =\int_0^1w(a)(1-a)^2\dd a=J_w(1),
\end{equation}
so the bound is sharp. The same Tonelli argument as before
handles every $L^2$ profile. More precisely the difference
$c(w)-J_w(u)$ is the sum of
$\sum_{n\ge1}(-\widehat K(n))\abs{\widehat F(n)}^2$ and
the overlap-variance term in~\eqref{eq:continuous-deficit}
with $a$ replaced by $w(a)$.
The use of nonpositive nonconstant Fourier coefficients is the
classical energy principle for a kernel of conditionally negative
type; see~\cite{BCR}. The preceding calculation states the
precise version needed here and proves it directly.

The symmetry condition is not arbitrary. For a bounded real
mean-zero $v$, differentiate at $\eta=0$ the profile
$(1+\eta v)/\norm{1+\eta v}_2$. The first variation is
\[
 2\int_0^1v(x)\left(
       \int_0^xK(a)\dd a+\int_0^{1-x}K(a)\dd a\right)\dd x.
\]
It vanishes for every such $v$ exactly when the bracket is
constant. Its derivative is $K(x)-K(1-x)$, so flat
stationarity is equivalent to kernel symmetry. Imaginary
first variations vanish automatically. The Fourier-sign
condition is a sufficient global-optimality condition beyond
stationarity; we do not assert that it is necessary.

There is also a finite-grid estimate without an asymptotic
optimizer argument. For a unit vector in $\C^d$, set
$F_i=\abs{v_i}^2$ as before. Then
\begin{align}\label{eq:taper-grid}
 \frac1d\sum_{j=1}^{d-1}w(j/d)\abs{c_j(v)}^2
 &\le\frac12\sum_{i,k=0}^{d-1}K((i-k)/d)F_iF_k\notag\\
 &=c(w)+\frac12\sum_{n\ne0}\widehat K(n)
              \left|\sum_iF_i e^{2\pi\ii ni/d}\right|^2
 \le c(w).
\end{align}
The kernel is periodic. All nonzero Fourier indices, including
multiples of $d$, are retained; their signs help the bound.
If $w$ is Lipschitz and $d=q+2M$, changing $j/d$ to $j/q$
and then enlarging the nonnegative sum proves
\begin{equation}\label{eq:taper-mesh-change}
 \frac1q\sum_{j=1}^{q-1}w(j/q)\abs{c_j(v)}^2
 \le\frac dq c(w)+\operatorname{Lip}(w)\frac dq
                                      \left(\frac dq-1\right).
\end{equation}
Here $\sum_j\abs{c_j(v)}^2\le d$ suffices for the error.

\subsection*{The full arithmetic deformation}

Let $g\in C^1([0,1];\R)$ be fixed, and suppose the profile
$w(a)=ag(a)^2$ satisfies the criterion above with $c(w)>0$.
Absolute Fourier convergence is automatic in this class:
the continuous periodic kernel has a weak derivative in $L^2$,
so Cauchy--Schwarz applied to its weighted Fourier coefficients
gives $\sum_n\abs{\widehat K(n)}<\infty$.
Define
\begin{equation}\label{eq:taper-operator}
 W_{f,\ell}^g=\sum_{p^m<e^\ell}b_f(p^m)g(m\log p/\ell)
                           (S_{m\log p}+S_{-m\log p}).
\end{equation}
For fixed $f,g,C$ and $2\le\abs t\le C\ell$, the complete
uniform conclusion is
\begin{equation}\label{eq:taper-pressure}
 c(w)t^2-O_{f,g,C}(1)
 \le\log\tau_\ell(e^{tW_{f,\ell}^g/\ell})
 \le c(w)t^2+O_{f,g,C}\bigl(1+\abs t\sqrt{\log(2+\abs t)}\bigr).
\end{equation}
We detail the changes to the pressure proof to keep every cost
explicit. The local absolute and squared budgets are multiplied
by at most $\norm g_\infty$ and $\norm g_\infty^2$,
respectively. The relative Taylor argument therefore still
restores all local harmonics at bounded cost for bounded
physical coupling; each prime retains its single phase.
The lower Bessel bound uses
$F(a)=g(a)^2(1-a)^2$. Stieltjes integration by parts against
\eqref{eq:variance}, using $F(1)=0$, gives
\[
 \sum_pb_p^2F(h_p/\ell)=\ell^2\int_0^1aF(a)\dd a+O_{f,g}(1)
                         =c(w)\ell^2+O_{f,g}(1).
\]
The fourth-power remainder remains summable.

For the upper bound take $M=\lceil D_gt^2\rceil$ and
$q=\lceil t^6\rceil$, choosing $D_g$ large enough to
absorb the factor $\norm g_\infty^2$ in the unmasked
word-tail estimate. The mask domination and relative Gaussian
replacement are unchanged. In the normalized weak-variance sum,
replacing $g(h_p/\ell)^2$ by $g(j_p/q)^2$ costs at most
\[
 \frac{\operatorname{Lip}(g^2)}q\frac{V_f(\ell)}{\ell^2}
                              =O_{f,g}(q^{-1}).
\]
In each long bin the main mass is now
$w(j/q)/q+g(j/q)^2/(2q^2)$; the summed second term is
$O_g(d/q^2)$. The short-bin budget remains
$O_{f,g}(\log^4(2+\ell)/\ell^2+q^{-2})$.
The moving-bin errors are unchanged up to a bounded factor.
Equation~\eqref{eq:taper-mesh-change} therefore bounds the
weak-variance half by
\[
 c(w)+O_{f,g,C}\left(M/q+q^{-1}
                              +\log^4(2+\ell)/\ell^2\right).
\]
The strong variance stays bounded. Gaussian concentration,
the same uniform absorption of the arithmetic error, and
restoration of the local strings prove~\eqref{eq:taper-pressure}.
The constants are uniform on any $C^1$-bounded family of
tapers satisfying the stated kernel conditions.

The proof of Section~\ref{sec:deviations} now gives rate
$x^2/(4c(w))$ for every diverging scale $r_\ell=O(\ell)$
of $W_{f,\ell}^g/(\ell r_\ell)$, with speed $r_\ell^2$.
This uses $c(w)>0$. It is not a central-law assertion obtained
by sending $r_\ell$ to a constant.

\subsection*{An explicit equal-variance family}

Take $g_\varepsilon(a)=\sqrt{1+\varepsilon\cos(2\pi a)}$.
The kernel is
$K_\varepsilon(a)=a(1-a)(1+\varepsilon\cos(2\pi a))$.
Multiplying its absolutely convergent Bernoulli series gives
\begin{align}\label{eq:taper-fourier}
 \widehat K_\varepsilon(0)&=\frac16-\frac{\varepsilon}{2\pi^2},
 &\widehat K_\varepsilon(1)&=-\frac1{2\pi^2}
                       +\varepsilon\left(\frac1{12}-\frac1{16\pi^2}\right),
 \notag\\
 \widehat K_\varepsilon(n)&=-\frac1{2\pi^2n^2}
       -\frac{\varepsilon}{4\pi^2}
          \left(\frac1{(n-1)^2}+\frac1{(n+1)^2}\right),&&n\ge2.
\end{align}
The $n=1$ restriction is
$\varepsilon\le24/(4\pi^2-3)$. For $n\ge2$ it is
$\varepsilon\ge-(n^2-1)^2/(n^2(n^2+1))$.
The positive ratio $(n^2-1)^2/[n^2(n^2+1)]$ increases with $n\ge2$
(as a function of $y=n^2$ its derivative has numerator
$(3y+1)(y-1)$). Thus all nonconstant coefficients are
nonpositive exactly on
\begin{equation}\label{eq:taper-range}
 -\frac9{20}\le\varepsilon\le\frac{24}{4\pi^2-3},
 \qquad c_\varepsilon=\frac1{12}-\frac{\varepsilon}{4\pi^2}>0.
\end{equation}
This closed interval lies strictly inside $(-1,1)$, so these
tapers are uniformly smooth and bounded away from zero.
The pressure estimate~\eqref{eq:taper-pressure} holds uniformly
over the entire interval, including its endpoints. This is
the exact interval for the Fourier-sign certificate, not a
claim that the flat profile ceases to be optimal outside it.

On a compact subinterval of its interior, both
$w_\varepsilon(a)/a$ and
$2\pi^2n^2[-\widehat K_\varepsilon(n)]$ have uniform positive
lower bounds. The exact weighted deficit consequently dominates
a positive constant times $D(u)$. Strong stability modulo
affine phase follows from~\eqref{eq:stability-intro}, uniformly
on each such compact parameter set. No endpoint stability
constant is asserted here.

This family has equal total leading first-prime variance:
\[
 \int_0^1ag_\varepsilon(a)^2\dd a=\frac12,
 \qquad
 \sum_pb_p^2g_\varepsilon(h_p/\ell)^2=\frac{\ell^2}{2}+O_f(1),
\]
uniformly over~\eqref{eq:taper-range}. The latter follows by
the same bounded-variation integration against~\eqref{eq:variance}.
Thus the unmasked product still has fixed-coupling logarithm
$s^2\ell^2/2+O_{f,s}(1)$, whereas its killed logarithm has
leading coefficient $s^2c_\varepsilon$. The paired-word identity
gives
\begin{equation}\label{eq:taper-survival}
 \log\E F_\ell
   =-s^2\left(\frac5{12}+\frac{\varepsilon}{4\pi^2}\right)\ell^2
                   +O_{f,s}(\ell\sqrt{\log(2+\ell)}).
\end{equation}
The distribution of the paired words here uses the tapered
first-prime coefficients. Formula~\eqref{eq:taper-survival}
isolates a change in boundary confinement while preserving
the total leading scalar variance.

Finally replace $W$ by $W^{g_\varepsilon}$ in
\eqref{eq:gamma-identity}, retaining the same gamma operator.
Its fixed-window Weyl coefficient satisfies
\begin{equation}\label{eq:taper-K}
 \log K_{f,\ell}^{\varepsilon}
  =\left(\frac1{48}-\frac{\varepsilon}{16\pi^2}\right)\ell^2
        +\log\frac{2\ell}{\sqrt N}
        +O_f(\ell\sqrt{\log(2+\ell)}),
\end{equation}
again with bounded lower error and constants uniform in the
displayed parameter interval. These are deformed logarithmic
operators. Only $\varepsilon=0$ is the original arithmetic
operator of the preceding sections.

\section*{Note on AI Use}
The author declares the use of consumer-grade Gemini-3 and GPT-6 models for prior-art search, exploring and developing ideas and proofs, drafting and LaTeX typesetting; the author takes full and sole responsibility for all contents of this work.

\begin{thebibliography}{99}

\bibitem{AL}
A.~O.~L. Atkin and J. Lehner,
\emph{Hecke operators on $\Gamma_0(m)$},
Math. Ann. \textbf{185} (1970), 134--160.
\href{https://doi.org/10.1007/BF01359701}{doi:10.1007/BF01359701}.

\bibitem{BCR}
C. Berg, J.~P.~R. Christensen, and P. Ressel,
\emph{Harmonic Analysis on Semigroups: Theory of Positive Definite
and Related Functions},
Graduate Texts in Mathematics, vol.~100, Springer, New York, 1984.
\href{https://doi.org/10.1007/978-1-4612-1128-0}{doi:10.1007/978-1-4612-1128-0}.

\bibitem{BHS}
A. Bose, R.~S. Hazra, and K. Saha,
\emph{Convergence of joint moments for independent random patterned matrices},
Ann. Probab. \textbf{39} (2011), no.~4, 1607--1620.
\href{https://doi.org/10.1214/10-AOP597}{doi:10.1214/10-AOP597};
\href{https://arxiv.org/abs/1211.3864}{arXiv:1211.3864}.

\bibitem{BSS}
A. Bose, K. Saha, and P. Sen,
\emph{Some patterned matrices with independent entries},
Random Matrices Theory Appl. \textbf{10} (2021), article 2150030.
\href{https://doi.org/10.1142/S2010326321500301}{doi:10.1142/S2010326321500301}.
Updated author version:
\href{https://arxiv.org/abs/2203.05837v1}{arXiv:2203.05837v1} (2022).

\bibitem{BDJ}
W. Bryc, A. Dembo, and T. Jiang,
\emph{Spectral measure of large random Hankel, Markov and Toeplitz matrices},
Ann. Probab. \textbf{34} (2006), no.~1, 1--38.
\href{https://doi.org/10.1214/009117905000000495}{doi:10.1214/009117905000000495}.

\bibitem{ChandruB}
R. Arun Chandru,
\emph{Arithmetic identification from interval Weil spectra},
preprint, 2026.
\href{https://doi.org/10.5281/zenodo.22834353}{doi:10.5281/zenodo.22834353}.


\bibitem{ChandruW}
R. Arun Chandru,
\emph{Sharp Weyl laws and resonance transitions for logarithmic
operators with killed translations},
preprint, 2026.
\href{https://doi.org/10.5281/zenodo.22881457}{doi:10.5281/zenodo.22881457}.


\bibitem{Christ}
M. Christ,
\emph{A sharpened Hausdorff--Young inequality},
preprint, 2014.
\href{https://arxiv.org/abs/1406.1210}{arXiv:1406.1210}.

\bibitem{Deligne}
P. Deligne,
\emph{La conjecture de Weil. I},
Publ. Math. Inst. Hautes \'Etudes Sci. \textbf{43} (1974), 273--307.
\href{https://doi.org/10.1007/BF02684373}{doi:10.1007/BF02684373}.

\bibitem{GRR}
A.~M. Garsia, E. Rodemich, and H. Rumsey, Jr.,
\emph{On some extremal positive definite functions},
J. Math. Mech. \textbf{18} (1969), no.~9, 805--834.

\bibitem{GJ}
T. Gillespie and G. Ji,
\emph{A prime number theorem for Rankin--Selberg $L$-functions
over number fields},
preprint, 2009.
\href{https://arxiv.org/abs/0910.3660}{arXiv:0910.3660}.

\bibitem{Gross}
L. Gross,
\emph{Logarithmic Sobolev inequalities},
Amer. J. Math. \textbf{97} (1975), no.~4, 1061--1083.
\href{https://doi.org/10.2307/2373688}{doi:10.2307/2373688}.

\bibitem{LW}
A. Laptev and T. Weth,
\emph{Spectral properties of the logarithmic Laplacian},
Anal. Math. Phys. \textbf{11} (2021), article 133.
\href{https://doi.org/10.1007/s13324-021-00527-y}{doi:10.1007/s13324-021-00527-y};
\href{https://arxiv.org/abs/2009.03395}{arXiv:2009.03395}.

\bibitem{Ledoux}
M. Ledoux,
\emph{The Concentration of Measure Phenomenon},
Mathematical Surveys and Monographs, vol.~89,
American Mathematical Society, Providence, RI, 2001.

\bibitem{Manov}
A.~D. Manov,
\emph{An extremal problem for positive definite functions with support in a ball},
Sb. Math. \textbf{215} (2024), no.~7, 920--931.
\href{https://doi.org/10.4213/sm10006e}{doi:10.4213/sm10006e}.

\bibitem{DLMF}
F.~W.~J. Olver, A.~B. Olde Daalhuis, D.~W. Lozier, et al., eds.,
\emph{NIST Digital Library of Mathematical Functions},
Section 5.11, asymptotic expansions of the gamma and digamma functions.
\url{https://dlmf.nist.gov/5.11}.

\bibitem{Onatski}
A. Onatski,
\emph{Precise asymptotics for the norm of large random rectangular Toeplitz matrices},
preprint, 2025.
\href{https://arxiv.org/abs/2509.04349}{arXiv:2509.04349}.

\bibitem{SV}
A. Sen and B. Vir\'ag,
\emph{The top eigenvalue of the random Toeplitz matrix and the sine kernel},
Ann. Probab. \textbf{41} (2013), no.~6, 4050--4079.
\href{https://doi.org/10.1214/13-AOP863}{doi:10.1214/13-AOP863};
\href{https://arxiv.org/abs/1109.5494}{arXiv:1109.5494}.

\bibitem{Tropp}
J.~A. Tropp,
\emph{User-friendly tail bounds for sums of random matrices},
Found. Comput. Math. \textbf{12} (2012), no.~4, 389--434.
\href{https://doi.org/10.1007/s10208-011-9099-z}{doi:10.1007/s10208-011-9099-z}.

\end{thebibliography}
\end{document}
